\section{Problem Formulation}
\label{sec:formulation}

To establish a rigorous foundation for \textbf{agent self-evolution}, we formalize the learning problem as enabling an LLM-based agent to autonomously improve its policy through interaction with the environment---\textit{without relying on predefined task distributions or reward functions}. 
Our approach conceptually separates this problem into two components: (1) an interaction sandbox $\mathcal{E}$, which defines the environment dynamics (states, actions) but provides no objectives, and (2) an unknown target task distribution $p_{\text{target}}(g)$, representing the `golden' tasks the agent must ultimately master. The core challenge is thus for the agent to leverage open-ended exploration in $\mathcal{E}$ to effectively estimate $p_{\text{target}}(g)$, and in turn, autonomously generate its own proxy training objectives (tasks) and learning signals (rewards) to guide continual improvement. We now formally define these components.


\paragraph{Environment Formulation}
{
To distinguish our setting from standard RL environments, 
we define an \textit{interaction sandbox} that specifies the agent’s accessible state space, 
action space, and transition dynamics, but does not provide any predefined reward function or task objective. 
Formally, the interaction sandbox is represented as:
\begin{equation}
\label{eq:interaction_world}
\mathcal{E} = (\mathcal{S}, \mathcal{A}, \mathcal{P}),
\end{equation}
where $\mathcal{S}$ denotes the set of observable states, 
$\mathcal{A}$ the set of executable actions, 
and $\mathcal{P}(s' \mid s, a)$ the transition probability distribution governing the environment dynamics. 
Unlike a standard MDP environment ${E} = (\mathcal{S}, \mathcal{A}, \mathcal{P}, r, \gamma)$, 
the interaction sandbox $\mathcal{E}$ excludes the reward function $r$ as well as the discount factor $\gamma$, 
reflecting an open-ended, non-rewarding setting in which the agent must autonomously generate 
its own learning signals and objectives to drive continual policy improvement.}

\paragraph{Target Task Space and Oracle Reward}

The agent is ultimately evaluated on a target task distribution $p_{\text{target}}(g)$ over tasks $g \in \mathcal{G}$, where $\mathcal{G}$ denotes the task space. Each task $g$ is associated with a desired goal state $s_g \in \mathcal{S}$, representing a desired terminal state or outcome (e.g., a correct calculation, a retrieved fact, or a completed workflow). Intuitively, executing a task $g$ means steering the environment from an initial state toward its corresponding goal state $s_g$.

For each task $g$, a ground-truth reward function $R_g(s, a)$ measures the utility of taking action $a$ in state $s$ with respect to achieving the goal $s_g$. Starting from an initial state $s_0 \sim p_0$, the agent aims to learn a goal-conditioned policy $\pi_\theta(a \mid s, g)$ that maximizes the expected return across the target task distribution:
\begin{equation}
\label{eq:target_objective}
J_{\text{target}}(\theta) = 
\mathbb{E}_{g \sim p_{\text{target}},\, s_0 \sim p_0}
\left[ 
V^{\pi_\theta}(s_0, g)
\right],
\end{equation}
where the goal-conditioned value function follows the standard formulation \citep{schaul2015universal, liu2022goal}:
\begin{equation}
\label{eq:value_function}
V^{\pi_\theta}(s_0, g) 
= 
\mathbb{E} 
\left[ 
\sum_{t=0}^{\infty} \gamma^t R_g(s_t, a_t) 
\,\middle|\, 
s_0, g, \pi_\theta 
\right],
\end{equation}
which generalizes value estimation across both environmental states and goal-conditioned tasks.


\paragraph{Proxy Objective of Self-Evolution}
Since the target task distribution and reward functions are unknown a priori, the agent must synthesize its own training tasks and reward signals through formal generation mechanisms:

\begin{itemize}
    \item \textbf{Proxy Task Generation:} The agent explores the environment via a \texttt{self-questioning} mechanism (Section~\ref{sec:self-questioning}) to produce a candidate set of solvable tasks $\mathcal{G}_{\text{candidate}}$. This process can be formalized as a mapping from the environment to a task distribution:
    \begin{equation}
    \label{eq:task_mapping}
    F_{\text{task}}: \mathcal{E} \to \Delta(\mathcal{G}), \quad \text{where} \quad p_{\text{train}} = F_{\text{task}}(\mathcal{E})
    \end{equation}
    The resulting distribution $p_{\text{train}}(g)$ emphasizes tasks of appropriate difficulty and diversity, emulating the unknown $p_{\text{target}}(g)$.

    \item \textbf{Proxy Reward Design:} The agent infers a proxy reward function through \texttt{self-attributing} (Section~\ref{sec:self-attribution}), formalized as a mapping from environment and tasks to reward functions:
    \begin{equation}
    \label{eq:reward_mapping}
    F_{\text{reward}}: \mathcal{E} \times \mathcal{G} \to (\mathcal{S} \times \mathcal{A} \to \mathbb{R}), \quad \text{where} \quad \hat{R}_g = F_{\text{reward}}(\mathcal{E}, g)
    \end{equation}
    This enables finer-grained credit assignment than sparse environment rewards, allowing the agent to learn from richer feedback without human intervention.
\end{itemize}

The core objective of self-evolution is to formulate and optimize $F_{\text{task}}$ and $F_{\text{reward}}$ functions such that maximizing the proxy training objective:
\begin{equation}
\label{eq:train_objective}
J_{\text{train}}(\theta) = \mathbb{E}_{g \sim F_{\text{task}}(\mathcal{E})} \left[ \hat{V}^{\pi_\theta}(s_0, g) \right],
\end{equation}
with 
\begin{equation}
\label{eq:proxy_value}
\hat{V}^{\pi_\theta}(s_0, g) = \mathbb{E} \left[ \sum_{t=0}^{\infty} \gamma^t F_{\text{reward}}(\mathcal{E}, g)(s_t, a_t) \,\middle|\, s_0, g, \pi_\theta \right]
\end{equation}
leads to the improvement of the true target objective in Eq. \ref{eq:target_objective}. Moreover, it optimizes exploration through a self-navigating mechanism, guiding $\pi_\theta$ to generate high-learning-value trajectories and addressing efficiency bottlenecks in multi-turn environments.

In other words, this manuscript aims to formulate $F_{\text{task}}$ and $F_{\text{reward}}$ through systematic methodological development such that the policy $\pi_\theta$ optimized under the proxy tasks and rewards performs well on the true target task distribution. This formulation establishes a systematic framework where the agent's self-improvement is guided by carefully formulated task and reward generation mechanisms as well as efficient learning methods.